\chapter{Introducción}
\label{cap:Introduccion}

La codificación clínica es un proceso esencial en la gestión sanitaria: traduce los diagnósticos a códigos estandarizados y sirve de puente entre la práctica clínica, la administración hospitalaria y la investigación médica.

Este trabajo se centra en la codificación automática mediante modelos BERT multilabel para CIE-10, un problema en el que confluyen la creciente complejidad de los sistemas de clasificación médica, los desafíos operativos de los entornos clínicos reales y las oportunidades que abre el Procesamiento del Lenguaje Natural (PLN).

\section{Motivación}

La motivación de este trabajo surge de varios problemas concretos:
\begin{itemize}
\item La correcta clasificación requiere conocimiento médico especializado, lo que limita la escalabilidad del proceso manual.
\item La normativa CIE-10 contiene 68.069 códigos posibles, lo que dificulta la tarea de clasificación y eleva el número de errores en el proceso manual.
\item Desde el 1 de enero de 2016 la CIE-10-ES es la clasificación de referencia para la codificación clínica en España, como indica el Real Decreto 69/2015~\cite{rd692015}, aunque su adopción obligatoria se ha extendido de forma gradual.
\item El tiempo medio de codificación de un alta hospitalaria es de 14 minutos, diez más que con la CIE-9, tal y como se indica en el informe de implantación de la CIE-10~\cite{sanidad2016cie10}.
\end{itemize}

La solución propuesta combina técnicas de \emph{Deep Learning} con conocimiento clínico estructurado para automatizar el proceso de codificación, reservando la revisión humana para los casos de baja certeza. El aprendizaje automático aplica las normas CIE-10 de forma consistente y acompaña cada código propuesto de una puntuación de confianza, lo que permite al codificador concentrarse en los casos ambiguos.


\section{Contexto disciplinar y tecnológico}

El análisis de trabajos previos revela tres generaciones tecnológicas:

\textbf{Sistemas basados en reglas} (1990--2010): Utilizan ontologías médicas como SNOMED-CT~\cite{stearns2001snomedct} con lógica booleana. Están muy limitados por su alta rigidez ante la variabilidad lingüística.

\textbf{Modelos estadísticos} (2010--2018): Aplicación de SVM y \emph{Random Forests} sobre características léxico-sintácticas. Alcanzan resultados competitivos en la CIE-9 pero no consiguen adaptarse a la mayor granularidad de la CIE-10~\cite{perotte2014}.

\textbf{Deep Learning} (2018-actual): Modelos basados en transformers con atención multicabezal, que desplazaron a las arquitecturas BiLSTM previas. Los resultados de CodiEsp ilustran bien el estado del arte: el mejor sistema por MAP\footnote{\emph{Mean Average Precision} (MAP): métrica oficial del \emph{benchmark} CodiEsp. Evalúa la calidad del \emph{ranking} de códigos predichos sin depender de un umbral de decisión fijo; a diferencia del F1-score, premia que el código correcto aparezca en las primeras posiciones de la lista.} fue IXA-AAA~\cite{miranda2020}, que combinó un clasificador XGBoost por código con BERT Multilingüe y alcanzó MAP~0,593. Los sistemas basados exclusivamente en transformers, BETO (\emph{The Mental Strokers}) y BERT-SciELO (ICB-UMA~\cite{lopezgarcia2020icb}), obtuvieron MAP~0,517 y 0,482 respectivamente.

Un caso ilustrativo en el contexto de CodiEsp es el del equipo IAM del \emph{Service d'Information Médicale} de Burdeos, que obtuvo el mayor F1-score de la tarea de 2020 (0,687) sin recurrir al Deep Learning~\cite{cossin2020}. Su sistema se basa íntegramente en un diccionario con estructura de árbol, construido a partir de las anotaciones de CodiEsp y terminología del CIE-10. La detección de entidades se realiza por coincidencia exacta, distancia de Levenshtein y un diccionario de abreviaturas; los términos que generaban muchos falsos positivos se eliminaron manualmente. Su precisión de 0,82 en diagnósticos, procesando todo el corpus en menos de 5 segundos en un portátil, sitúa al enfoque por diccionario por encima de todos los sistemas neuronales en esta métrica. Sin embargo, el MAP del equipo IAM fue de 0,521, inferior al 0,593 de IXA-AAA: la coincidencia exacta produce predicciones de alta precisión pero un \emph{ranking} menos fino que los modelos probabilísticos.

Este contraste ilustra una compensación conocida: los enfoques por diccionario tienden a superar a los modelos neuronales en precisión cuando el corpus de entrenamiento es reducido, como ocurre en CodiEsp, ya que si el término está en el léxico, la asignación es directa y muy fiable. Sin embargo, son intrínsecamente rígidos ante formulaciones no previstas o variantes lingüísticas no indexadas. El presente trabajo incluye un clasificador de diccionario como \emph{baseline} de referencia, pero el foco es el clasificador neuronal basado en RigoBERTa-Clinical, cuyo valor añadido reside en la capacidad de generalizar a expresiones no vistas en el diccionario.

\subsection{Avances en transformers médicos}
Los modelos BERT adaptados al dominio clínico, como ClinicalBERT~\cite{alsentzer2019} y BioBERT~\cite{lee2020}, destacan por capturar las relaciones contextuales de las narrativas médicas y por manejar la ambigüedad terminológica: el término \emph{EPOC}, por ejemplo, puede designar el diagnóstico principal (código J44.x) o una comorbilidad que condiciona el tratamiento, con implicaciones distintas en la codificación. Su capacidad de transferencia de aprendizaje entre especialidades médicas los hace especialmente adecuados para este dominio.

Este trabajo extiende estos enfoques mediante:
\begin{enumerate}
\item Adaptación multilingüe para textos clínicos en español mediante RigoBERTa-Clinical.
\item Clasificador multilabel plano sobre el espacio completo de códigos CIE-10.
\item Función de pérdida BCE ponderada por frecuencia inversa de clase, con límite de peso por clase para mitigar el desequilibrio extremo.
\item Preentrenamiento débil sobre las descripciones oficiales del catálogo CIE-10-ES.
\end{enumerate}

\subsection{Retos pendientes}
La literatura identifica cuatro desafíos no resueltos:
\begin{itemize}
\item \textbf{Clasificación multietiqueta a gran escala}: la CIE-10 abarca decenas de miles de códigos y cada caso recibe varios simultáneamente: en el corpus CodiEsp~\cite{miranda2020}, una media de 11,3 códigos de diagnóstico por informe. Esto exige conjuntos de datos muy amplios para clasificar de forma consistente.
\item \textbf{Distribución de los códigos en la codificación}: La distribución de los códigos es muy asimétrica: hay familias de códigos muy frecuentes y otras con presencia casi nula, como las enfermedades raras.
\item \textbf{Evaluación clínica}: Las métricas tradicionales (F1, precisión) no capturan las consecuencias clínicas de los errores ni el grado de relación entre el código predicho y el correcto, lo que limita la capacidad de diseñar funciones de pérdida adaptadas al dominio.
\item \textbf{Explicabilidad}: El Reglamento UE 2017/745~\cite{eu2017745} clasifica el software de apoyo a la decisión clínica como producto sanitario de clase IIa o superior cuando influye en decisiones diagnósticas, exigiendo documentación de la lógica de decisión, trazabilidad de las predicciones y capacidad de auditoría. Para un clasificador automático de códigos CIE-10 esto implica que cada predicción debe acompañarse de una puntuación de confianza interpretable y un registro inmutable del modelo y la versión empleados.
\end{itemize}

\section{Estructura de la memoria}

El resto del documento se organiza como sigue. El capítulo~\ref{cap:Objetivo} fija el objetivo del trabajo y lo descompone en subobjetivos. El capítulo~\ref{cap:Planificacion} describe la metodología de trabajo, los recursos empleados, la planificación temporal y la estimación de costes y riesgos. El capítulo~\ref{cap:Desarrollo} constituye el núcleo técnico: establece los requisitos funcionales y no funcionales que debe cumplir el sistema y recoge los modelos de análisis, la obtención y el análisis de los datos, el diseño y el proceso de entrenamiento del clasificador, la arquitectura del sistema completo (motor de IA, backend, frontend e infraestructura) y las prácticas de calidad aplicadas. El capítulo~\ref{cap:Conclusiones} valora el grado de cumplimiento de los objetivos y los requisitos, expone las líneas de trabajo futuro y cierra con la valoración personal.

La memoria se completa con una serie de anexos que amplían aspectos concretos sin interrumpir la lectura del cuerpo principal: la comparativa de modelos candidatos (Anexo~\ref{cap:AnexoA}), la estructura del repositorio y el entorno de ejecución (anexos~\ref{cap:AnexoB} y~\ref{cap:AnexoK}), la especificación de las interfaces entre servicios (anexos~\ref{cap:AnexoC} y~\ref{cap:AnexoM}), el modelo de datos (Anexo~\ref{cap:AnexoD}), las decisiones de arquitectura (Anexo~\ref{cap:AnexoE}), el proceso de aprendizaje y la configuración final del modelo (anexos~\ref{cap:AnexoF} y~\ref{cap:AnexoH}), el análisis estadístico del corpus (Anexo~\ref{cap:AnexoG}), la arquitectura de producción (Anexo~\ref{cap:AnexoI}), el tablero de tareas (Anexo~\ref{cap:AnexoJ}) y la aumentación de datos mediante retrotraducción (Anexo~\ref{cap:AnexoL}).
